% Round 15. Multiplicity/native extension and explicit transpose criterion.
\section{A numerical component condition for higher-dimensional codimension two}
\label{r15model:section}

The general degree-two modification argument does not supply a
numerically admissible lift on a semistable model.  The older curve
proof controls a different transpose kernel in dimensions greater
than three.  We isolate that distinction and give a sufficient
component condition under which the missing kernel is nevertheless
controlled by the already established divisor argument.

\paragraph{Historical provenance.}
The reduced strictly semistable primitive normalization, isotropy,
and positivity mechanism already occurs in Round 7,
\path{history_01_07/bb_round7/proofs/threefold_codimension_two_obstruction.tex}
and \path{threefold_boundary_primitive_obstruction.tex}.
Those sources already assume geometric rational homological
vanishing.  The present section records the component criterion
with arbitrary scheme multiplicities, its exact transpose kernel,
native residue degrees, and the revised pure-operational actual
bridge of Round 12.  Its combination with the new analytic dominant
descent is a Round 15 consequence; the reduced positivity mechanism
is not claimed as a new discovery.

Fix a finite valued extension $K/\mathbf Q_p$, with valuation ring
$R$ and finite residue field $k$.  Let $\mathcal X/R$ be an integral
regular flat projective model of a smooth geometrically integral
projective $n$-fold $X$, where $n\geq3$.  Suppose
\[
 \mathcal X_k=\sum_{i\in I}m_iD_i,\qquad m_i>0,
\tag{LW15.1}\label{r15model:snc}
\]
has strict normal crossings reduced support with all closed strata
smooth over $k$.  Fix a relatively ample line bundle $H$ on
$\mathcal X$; write $L$ for intersection with its restriction on
any stratum.  All numerical groups, cycle groups, and degrees below
are rational and native over $k$.

\begin{definition}[The degree-two component condition]
\label{r15model:condition}
The pair $(\mathcal X,H)$ satisfies condition $\mathrm{LC}_2$ if,
for every component $D_i$, the map
\[
 L:N^1(D_i)_{\mathbf Q}\longrightarrow N^2(D_i)_{\mathbf Q}
\quad\hbox{is surjective}.
\tag{LW15.2}\label{r15model:LC2}
\]
The same $H$ must be used on every component and every intersection.
This is a hypothesis on numerical groups, not a claim that arbitrary
semistable alterations possess it.  It does not require a
cycle-class map out of $N^2$ or a higher-codimension positivity
conjecture.
\end{definition}

\begin{lemma}[Rank form of the component condition]
\label{r15model:rank-condition}
For a smooth integral projective $n$-fold $D$ over the residue
field, $n\geq3$, the condition $L_H:N^1(D)\to N^2(D)$ being
surjective is equivalent to
\[
                 \dim_{\mathbf Q}N^1(D)=\dim_{\mathbf Q}N^2(D).
\]
Consequently its truth is independent of the choice of ample $H$.
The powers used to compare different strata must still come from
one global relatively ample class in the model argument.
\end{lemma}

\begin{proof}
For every ample $H$, the divisor Hodge index pairing
$\deg_D(u vH^{n-2})$ is nondegenerate.  If $H u=0$ in $N^2(D)$,
all these pairings vanish and $u=0$.  Thus $L_H$ is always
injective.  It is surjective exactly when the two finite dimensions
are equal.  This proves both assertions, without a higher-degree
hard Lefschetz or positivity assertion.
\end{proof}

\begin{lemma}[Two numerical Lefschetz isomorphisms]
\label{r15model:lefschetz}
Let $D$ be a smooth integral projective $n$-fold over $k$, $n\geq3$,
and let $H$ be ample.  If $L:N^1(D)\to N^2(D)$ is surjective, then
\[
 L^{n-3}:N^1(D)\xrightarrow{\sim}N^{n-2}(D),\qquad
 L^{n-3}:N^2(D)\xrightarrow{\sim}N^{n-1}(D).
\tag{LW15.3}\label{r15model:two-isomorphisms}
\]
\end{lemma}

\begin{proof}
The divisor Hodge index theorem used in frozen Round 12 makes
\[
 (u,v)\longmapsto\deg_D(u v H^{n-2})
\]
nondegenerate on $N^1(D)_{\mathbf Q}$.  Therefore $L:N^1\to N^2$
is injective: if $Lu=0$, all these pairings vanish.  By the stated
surjectivity it is an isomorphism.  Numerical intersection gives
perfect finite-dimensional pairings $N^a(D)\times N^{n-a}(D)$.
Hence $\dim N^{n-2}=\dim N^2=\dim N^1=\dim N^{n-1}$.

If $L^{n-3}u=0$ in $N^{n-2}$, pairing with $Lv\in N^2$ gives
$\deg_D(u v H^{n-2})=0$ for all $v$, and $u=0$.
This proves the first isomorphism by the dimension equality.
For the second, write every $w\in N^2$ uniquely as $Lv$.
If $L^{n-3}w=0$, pairing with an arbitrary $u\in N^1$ gives
the same nondegenerate divisor pairing of $u$ with $v$.
Thus $v=w=0$, and the second map is an isomorphism as well.
For $n=3$ the asserted maps are identities; the component
condition is automatic by the same divisor index pairing.
\end{proof}

Set $V^a=\bigoplus_i N^a(D_i)_{\mathbf Q}$ and retain the
frozen actual vertical incidence maps
\[
 A_a:V^a\longrightarrow V^{a+1},\qquad
 A_a(v)=\left(j_i^!\sum_j j_{j*}v_j\right)_i,
\tag{LW15.4}\label{r15model:A}
\]
where $j_i:D_i\hookrightarrow\mathcal X$.
Their numerical definition is the restriction--Gysin formula on
double strata together with the self-normal term on each diagonal.
Let $M$ multiply the $i$th component by $m_i$.
Write $\delta$ for divisor restriction differences on double
strata and $\gamma_0$ for the signed Gysin of constants on those
strata.  These symbols have exactly the conventions of Round 12,
\path{numerical_lift/curve_lifts_and_actual_green.tex},
(NL3)--(NL8).  In particular the already proved formula is
\[
                  \ker A_1=M\ker\delta+\operatorname{im}\gamma_0.
\tag{LW15.5}\label{r15model:old-kernel}
\]
No unweighted version of this formula is used.

\begin{lemma}[The required transpose kernel under $\mathrm{LC}_2$]
\label{r15model:kernel}
Under condition $\mathrm{LC}_2$,
\[
 \ker A_{n-2}
   = L^{n-3}\ker A_1
   = L^{n-3}\bigl(M\ker\delta+\operatorname{im}\gamma_0\bigr).
\tag{LW15.6}\label{r15model:kernel-reduction}
\]
\end{lemma}

\begin{proof}
Restriction and Gysin commute with intersection by the globally
chosen $H$.  The normal self-intersection term also commutes with
it.  Thus
\[
             A_{n-2}L^{n-3}=L^{n-3}A_1.
\]
The first isomorphism of Lemma~\ref{r15model:lefschetz}, component
by component, writes every $v\in V^{n-2}$ uniquely as $L^{n-3}w$.
The second makes $L^{n-3}:V^2\to V^{n-1}$ injective.  Therefore
$A_{n-2}v=0$ if and only if $A_1w=0$.
This proves the first equality, and the frozen weighted divisor
kernel formula proves the second.  No assertion is made that a
numerical hard Lefschetz statement in higher degree holds without
$\mathrm{LC}_2$.
\end{proof}

\begin{theorem}[A numerically admissible codimension-two lift]
\label{r15model:lift}
Assume \eqref{r15model:snc} and $\mathrm{LC}_2$.
Let $z\in\operatorname{CH}^2(X)_{\mathbf Q}$ satisfy
\[
 \operatorname{cl}_\ell(z)=0
       \quad\hbox{in }H^4(X_{\bar K},\mathbf Q_\ell(2))
       \qquad(\ell\ne p).
\]
Every extension $\Theta\in\operatorname{CH}^2(\mathcal X)_{\mathbf Q}$
of $z$ admits a correction by vertical divisors on the components,
\[
 \Theta^{\mathrm{ad}}=\Theta-\sum_i j_{i*}u_i,
       \qquad u_i\in\operatorname{CH}^1(D_i)_{\mathbf Q},
\tag{LW15.7}\label{r15model:correction}
\]
such that
\[
         j_i^!\Theta^{\mathrm{ad}}=0
                         \quad\hbox{in }N^2(D_i)_{\mathbf Q}
                         \quad\hbox{for every }i.
\tag{LW15.8}\label{r15model:component-zero}
\]
There is a pure operational virtual bundle
$\xi\in K_0(\mathcal X)_{\mathbf Q}$ with generic
$\operatorname{ch}_2(\xi|_X)=z$ and degree zero on every
integral two-dimensional special-fiber subvariety.
Under the frozen model-Chern detector and exact-support bridge,
$\mathfrak o_X^2(z)=0$, and every prescribed rational cycle in
class $z$ has a strong G form over $K$, on its exact complement,
against every actual compact test.
\end{theorem}

\begin{proof}
Put $T_i=j_i^!\Theta$ in $N^2(D_i)$ and consider the rational
functional on $V^{n-2}$
\[
                    t(v)=\sum_i\deg_{D_i}(T_i v_i).
\]
We prove that it annihilates the two summands of
\eqref{r15model:kernel-reduction}.

First, the signed double-stratum Gysin identity
$\sum_i j_{i*}(\gamma_0 a)_i=0$ holds on the common support,
with the two occurrences of each double stratum cancelling.
Intersecting by the same $H^{n-3}$ and applying the projection
formula makes the two supported degree contributions from $T$
cancel as well.  Hence
\[
                         t(L^{n-3}\gamma_0a)=0.
\]
This computation takes degrees on the proper special strata.
It does not use an unweighted degree homomorphism on
$\operatorname{CH}_0(\mathcal X)$.

Next take $w\in\ker\delta$.  The frozen divisor cohomological-lift
lemma, Round 12 \texttt{r12n:cohom-lift}, gives
$b\in H^2(\mathcal X_{R^{\mathrm{sh}}},\mathbf Q_\ell(1))$
whose component restrictions are the divisor classes of $w_i$.
It uses numerical-to-homological triviality only for divisors,
the finite-field purity of the smooth strata, and proper base
change.  Let $a=\operatorname{ch}_2(\xi_\Theta)$ be the
cohomological class of the pure operational representative of
$\Theta$, as in Round 12 \texttt{r12n:K0}, and let
$\mathfrak h=c_1(H)$ in cohomology.  Apply the frozen weighted
trace lemma \texttt{r12n:weighted-trace} to the top-degree class
\[
                           a\cup b\cup\mathfrak h^{n-3}.
\]
It gives
\[
 \begin{split}
 t(L^{n-3}Mw)
 &=\sum_i m_i\deg_{D_i}(T_i w_i H^{n-3})\\
 &=\operatorname{Tr}_{X_{\bar K}}
       \bigl(\operatorname{cl}_\ell(z)\cup b|_{X_{\bar K}}
                                    \cup\mathfrak h^{n-3}|_{X_{\bar K}}\bigr)
 =0.
 \end{split}
\tag{LW15.9}\label{r15model:weighted-annihilation}
\]
The class $a$ comes from the chosen Chow representative, not from
an unproved cycle-class map on $N^2$.  The factor $b$ may be any
of the cohomological lifts supplied by the frozen lemma; no
Frobenius-invariant lift or Tate-conjecture surjectivity is needed.

Lemma~\ref{r15model:kernel} now gives $t|_{\ker A_{n-2}}=0$.
The perfect numerical pairings identify the transpose of
$A_{n-2}:V^{n-2}\to V^{n-1}$ with
$A_1:V^1\to V^2$; this is the same restriction--Gysin and normal
term calculation as in Round 12.  Finite-dimensional rational
linear algebra yields
\[
                         T\in(\ker A_{n-2})^\perp
                                          =\operatorname{im}A_1.
\]
Choose $u\in V^1$ with $A_1u=T$, and choose rational divisor
representatives on the smooth $D_i$.  Formula
\eqref{r15model:correction} has unchanged generic class and its
component restrictions are $T-A_1u=0$, proving
\eqref{r15model:component-zero}.

Apply the frozen pure-operational representative lemma
\texttt{r12n:K0} in degree $r=2$ to $\Theta^{\mathrm{ad}}$.
It explicitly applies in arbitrary degree, even though the main
curve theorem used it with $r=n-1$.
It produces $\xi$ with all other operational Chern-character
components zero.  Every integral surface in the special fiber
lies in some $D_i$.  Its supported degree equals the numerical
pairing of $j_i^!\Theta^{\mathrm{ad}}$ with that surface, hence
is zero.  This includes singular surfaces and surfaces in
multiple components.

The frozen numerical model-Chern detector makes
$\operatorname{ch}_2(\bar\xi)=0$ as an actual form.  The
Round 11 generic-Chow-zero primitive theorem identifies its
model-curvature class with $\mathfrak o_X^2(z)$, and thus that
class is zero.  Apply the established sufficient direction
$\mathfrak o=0\Rightarrow G$ to the prescribed representative.
It subtracts a global actual primitive from its exact-support
Green pair.  The correction adds no new pole and its identity
holds on all actual compact tests.  No reverse implication from
G is used.
\end{proof}

\begin{remark}[Readily checkable component groups]
\label{r15model:examples}
A component with $N^1=\mathbf Q[H]$ and
$N^2=\mathbf Q[H^2]$ satisfies $\mathrm{LC}_2$.
In particular projective-space components satisfy it by their
Chow ring.  This verifies a component condition when such a
regular SNC model is supplied; it is not an existence theorem
for a model with these components.

More generally, on a component for which the numerical groups
and ample class are given explicitly, the condition is exactly
surjectivity of one finite rational multiplication matrix
$L:N^1\to N^2$.  No assertion about every codimension-two cycle
being a complete intersection on the original generic fiber
is part of the condition.
\end{remark}

\paragraph{The remaining numerical obstruction without this condition.}
For an arbitrary supplied regular SNC model and any fixed extension
$\Theta$, a correction of the displayed type exists if and only if
\[
           t_\Theta|_{\ker A_{n-2}}=0.
\tag{LW15.10}\label{r15model:exact-numerical-gate}
\]
This follows from the same transpose pairing, with no extra
hypothesis.  The frozen curve theorem proves the corresponding
annihilation on $\ker A_1$, which is the correct transpose
kernel for codimension $n-1$.  For $n>3$ those are different
requirements.  Our new component condition supplies a genuine
identification of the two kernels by \eqref{r15model:kernel-reduction};
without it, that identification is not asserted.  In particular,
this section does not prove a numerically admissible lift for all
geometrically homologically trivial codimension-two cycles on
arbitrary higher-dimensional SNC models.
